% Widen OLEBFont's inherited Utopia word spacing by 20%, retaining
% its proportions. The multiplier is 6/5 = 1.2.
% PTXprint uses plain XeTeX; set TeX's glue dimensions after font selection
% so every face, size and OpenType feature combination gets the same policy.
\begingroup
\catcode`\@=11
% fontname yields category-12 characters. Match with the same categories;
% the appended sentinel also handles fonts that do not contain OLEBFont.
\edef\BibleSpacingDefine{%
  \noexpand\gdef\noexpand\BibleSpacingMatch##1\detokenize{OLEBFont}##2\noexpand\relax}
\BibleSpacingDefine{%
  \def\BibleSpacingSuffix{#2}%
  \ifx\BibleSpacingSuffix\empty\else
    \fontdimen2\font=\dimexpr\fontdimen2\font*6/5\relax
    \fontdimen3\font=\dimexpr\fontdimen3\font*6/5\relax
    \fontdimen4\font=\dimexpr\fontdimen4\font*6/5\relax
    % Font dimensions persist across groups; so must this initialization flag.
    \expandafter\gdef\csname BibleSpacing:\fontname\font\endcsname{}%
  \fi
}
\gdef\BibleWordSpacing{%
  \ifcsname BibleSpacing:\fontname\font\endcsname\else
  \begingroup
  \edef\BibleSpacingName{\fontname\font\detokenize{OLEBFont}}%
  \expandafter\BibleSpacingMatch\BibleSpacingName\relax
  \endgroup
  \fi
}
% Preserve the existing selector, including the optical-margin adapter.
\ifx\BibleSpacingSetFont\undefined\global\let\BibleSpacingSetFont\s@tfont\fi
\gdef\s@tfont#1#2{%
  \BibleSpacingSetFont{#1}{#2}%
  \BibleWordSpacing
}
\endgroup
